\chapter{TỔNG QUAN TÀI LIỆU}

\section{Giới thiệu về phân cụm (Clustering)}

Phân cụm là một kỹ thuật học máy không giám sát (unsupervised learning) \cite{bishop_ml} nhằm nhóm các đối tượng dữ liệu thành các cụm một cách sao cho:

\begin{itemize}
    \item Các đối tượng trong cùng một cụm có sự tương đồng cao nhất
    \item Các đối tượng ở các cụm khác nhau khác nhau nhất
\end{itemize}

Phân cụm được ứng dụng rộng rãi trong:
\begin{itemize}
    \item Phân khúc thị trường (market segmentation)
    \item Phát hiện cấu trúc dữ liệu
    \item Nén dữ liệu (data compression)
    \item Bộ lọc cộng tác (collaborative filtering)
    \item Phát hiện bất thường (anomaly detection)
\end{itemize}

\section{Các phương pháp phân cụm chính}

\subsection{Phân cụm dựa trên khoảng cách (Distance-based Clustering)}

\textbf{K-means:}

K-means \cite{kmeans_original} là một trong những thuật toán phân cụm phổ biến nhất. Phiên bản cải tiến K-means++ \cite{kmeans_plus_plus} giải quyết vấn đề khởi tạo ngẫu nhiên. Thuật toán hoạt động như sau:

\begin{enumerate}
    \item Chọn ngẫu nhiên K tâm cụm
    \item Gán mỗi điểm dữ liệu tới tâm cụm gần nhất
    \item Tính toán lại tâm cụm bằng trung bình các điểm trong cụm
    \item Lặp lại bước 2-3 cho đến khi hội tụ
\end{enumerate}

Hàm mục tiêu của K-means được định nghĩa là:

\begin{equation}
J = \sum_{k=1}^{K} \sum_{x \in C_k} \|x - \mu_k\|^2
\end{equation}

Trong đó:
\begin{itemize}
    \item $C_k$ là cụm thứ $k$
    \item $\mu_k$ là tâm của cụm thứ $k$
    \item $K$ là số lượng cụm
\end{itemize}

\textbf{Ưu điểm:}
\begin{itemize}
    \item Đơn giản, dễ hiểu và triển khai
    \item Nhanh chóng với dữ liệu lớn
    \item Dễ dàng kết hợp với bộ kỳ vọng-tối đa (EM)
\end{itemize}

\textbf{Nhược điểm:}
\begin{itemize}
    \item Phải xác định trước số cụm K
    \item Nhạy cảm với khởi tạo ban đầu
    \item Có thể hội tụ tới cực tiểu cục bộ
    \item Giả định các cụm có hình dạng gần như hình cầu
\end{itemize}

\subsection{Phân cụm phân cấp (Hierarchical Clustering)}

Phân cụm phân cấp tạo ra một cây phân cấp (dendrogram) của các cụm. Có hai phương pháp chính:

\begin{itemize}
    \item \textbf{Agglomerative (Bottom-up):} Bắt đầu với mỗi điểm là một cụm và liên kết chúng lại
    \item \textbf{Divisive (Top-down):} Bắt đầu với tất cả các điểm trong một cụm và chia chúng ra
\end{itemize}

Các phương pháp liên kết phổ biến:
\begin{itemize}
    \item \textbf{Single Linkage:} Khoảng cách nhỏ nhất giữa hai cụm
    \item \textbf{Complete Linkage:} Khoảng cách lớn nhất giữa hai cụm
    \item \textbf{Average Linkage:} Khoảng cách trung bình giữa hai cụm
    \item \textbf{Ward Linkage:} Giảm thiểu phương sai \cite{hierarchical_clustering} (được sử dụng trong báo cáo này)
\end{itemize}

\subsection{Phân cụm dựa mật độ (Density-based Clustering)}

\textbf{DBSCAN (Density-Based Spatial Clustering of Applications with Noise)}

DBSCAN \cite{dbscan} nhóm các điểm dựa trên mật độ cục bộ. Các điểm được phân loại thành:
\begin{itemize}
    \item \textbf{Điểm lõi (Core points):} Có ít nhất MinPts điểm khác trong bán kính $\epsilon$
    \item \textbf{Điểm biên (Border points):} Không phải điểm lõi nhưng là hàng xóm của điểm lõi
    \item \textbf{Điểm nhiễu (Noise points):} Không phải điểm lõi cũng không phải điểm biên
\end{itemize}

\textbf{Ưu điểm:}
\begin{itemize}
    \item Không cần xác định K trước
    \item Phát hiện hình dạng cụm tùy ý
    \item Phát hiện điểm bất thường một cách tự nhiên
\end{itemize}

\textbf{Nhược điểm:}
\begin{itemize}
    \item Nhạy cảm với các tham số $\epsilon$ và MinPts
    \item Khó xử lý dữ liệu có mật độ biến đổi
    \item Chi phí tính toán $O(n^2)$ trong trường hợp xấu nhất (hoặc $O(n \log n)$ nếu dùng spatial indexing)
\end{itemize}

\section{Xác định số cụm tối ưu}

\subsection{Phương pháp Elbow (Elbow Method)}

Phương pháp Elbow dựa trên quan sát rằng:
\begin{itemize}
    \item Khi $K$ tăng, inertia (tổng bình phương khoảng cách) giảm
    \item Có một điểm "gãy khúc" (elbow) nơi giảm chậm lại
    \item Điểm này chỉ ra K tối ưu
\end{itemize}

Công thức inertia:
\begin{equation}
\text{Inertia} = \sum_{i=1}^{n} \|x_i - \mu_{c_i}\|^2
\end{equation}

\subsection{Chỉ số Silhouette (Silhouette Score)}

Silhouette Score \cite{silhouette_score} đo lường mức độ tương tự của một điểm với cụm của nó so với các cụm khác. Các độ đo khác như Davies-Bouldin Index \cite{davies_bouldin} và Calinski-Harabasz Index \cite{calinski_harabasz} cũng được sử dụng để đánh giá chất lượng phân cụm \cite{clustering_evaluation}.

\begin{equation}
s(i) = \frac{b(i) - a(i)}{\max\{a(i), b(i)\}}
\end{equation}

Trong đó:
\begin{itemize}
    \item $a(i)$ = khoảng cách trung bình từ điểm $i$ tới các điểm khác trong cùng cụm
    \item $b(i)$ = khoảng cách trung bình nhỏ nhất từ điểm $i$ tới các điểm trong các cụm khác
    \item $s(i) \in [-1, 1]$: giá trị cao hơn tốt hơn
\end{itemize}

Silhouette Score trung bình từ 0.5 đến 1 chỉ ra phân cụm tốt.

\section{Giảm chiều dữ liệu: Principal Component Analysis (PCA)}

PCA \cite{pca_paper} là một kỹ thuật giảm chiều dữ liệu tuyến tính. Nó tìm các thành phần chính (principal components) có phương sai cao nhất.

\begin{itemize}
    \item \textbf{PC1:} Hướng có phương sai lớn nhất
    \item \textbf{PC2:} Hướng trực giao với PC1 có phương sai lớn thứ hai
\end{itemize}

Công thức phương sai được giải thích:
\begin{equation}
\text{Variance Explained} = \frac{\lambda_i}{\sum_{j=1}^{m} \lambda_j}
\end{equation}

Trong đó $\lambda_i$ là eigenvalue thứ $i$.

\textbf{Ưu điểm:}
\begin{itemize}
    \item Giảm số chiều mà vẫn giữ lại phần lớn thông tin
    \item Giúp trực quan hóa dữ liệu cao chiều
    \item Giảm overfitting
    \item Tăng tốc độ tính toán
\end{itemize}

\section{Phân cụm khách hàng trong ngành tài chính}

Phân cụm khách hàng thẻ tín dụng \cite{customer_segmentation_review, financial_segmentation} là ứng dụng quan trọng để quản lý rủi ro \cite{credit_risk} và tối ưu hóa chiến lược kinh doanh. Phân cụm khách hàng thẻ tín dụng giúp:

\begin{enumerate}
    \item \textbf{Phát triển sản phẩm:} Tạo các sản phẩm phù hợp cho mỗi phân khúc
    
    \item \textbf{Giá định:} Đề xuất giá hoặc lãi suất khác nhau dựa trên phân khúc
    
    \item \textbf{Quản lý rủi ro:} Xác định nhóm rủi ro cao để kiểm soát tốt hơn
    
    \item \textbf{Tiếp thị được cá nhân hóa:} Tối ưu hóa các chiến dịch tiếp thị cho từng phân khúc
    
    \item \textbf{Tăng giữ chân:} Phát hiện khách hàng có nguy cơ rời bỏ cao và đề xuất các ưu đãi
    
    \item \textbf{Phát hiện gian lận:} Xác định hành vi bất thường có thể là gian lận
\end{enumerate}

\section{Trình trạng công nghệ hiện tại}

Các công cụ phân tích phổ biến trong phân cụm dữ liệu:
\begin{itemize}
    \item \textbf{R:} Thư viện stats, tidyverse, cluster
    \item \textbf{Python:} scikit-learn \cite{scikit_learn}, scipy, pandas, numpy
    \item \textbf{SAS, SPSS:} Công cụ phân tích thương mại
    \item \textbf{Tableau, Power BI:} Trực quan hóa và business intelligence
    \item \textbf{Apache Spark MLlib:} Xử lý dữ liệu phân tán quy mô lớn
\end{itemize}

Trong báo cáo này, Python được lựa chọn sử dụng \cite{hands_on_ml} vì:
\begin{itemize}
    \item Cộng đồng lớn và hỗ trợ tốt
    \item Thư viện mạnh mẽ cho khai phá dữ liệu (scikit-learn, pandas)
    \item Mã nguồn mở và miễn phí
    \item Dễ học và triển khai nhanh
    \item Tích hợp tốt với Jupyter Notebook cho phân tích tương tác
\end{itemize}
